\begin{abstract}
IoT device logs are long, identifier-dense, and routinely copied into language-model
prompts. A question about an account names exactly the value that must not reach the model,
so the queried entity is also the entity under protection. Redaction removes the value and,
with it, the handle the reader needs; compression keeps the handle but drops the record that
carries the answer. We show that the two goals need not be traded against each other. In
these logs the record that answers the question and the record that stores the protected
value are the same block, so \textsc{DARK} keeps that block and aliases the values inside it.

The rule is training-free: it labels each unit on sensitivity and relevance, applies
one of four operations---delete, alias, refill, keep---and writes the same alias map
into the question. Relevance is read from the question alone, so the compressor never
consults the reference answer.

On $2{,}728$ cases at a $50\%$ budget \textsc{DARK} reaches F1 $0.4937$, level with the
strongest query-aware compressor ($0.4831$; difference not significant) while leaving $0\%$
of protected values recoverable against that method's $98.1\%$, and exceeds the only leak-free
baseline by $0.3226$ points. On a 112-case set, keeping $48.6\%$ of the tokens is
statistically indistinguishable from aliasing the full context without compression ($0.6366$
vs.\ $0.6475$). Removing the alias step costs $0.3273$ F1 where the anchor is itself a
protected value and $0.0013$ where the question names no entity, so the rule acts where designed
and is inert elsewhere. The ordering holds at every budget from $10\%$ to $75\%$,
across a second reader and two paraphrases.
\end{abstract}
